% Suggested insertion for Geometry Before Fields
% Status: cosmology-owned R1 result; no BHSM microscopic normalization claim.
% Recommended placement: after the action-native matter/backreaction section
% and before observable-specific kernels.

\section{Coupled environmental realization of the physical $n=2$ state}
\label{sec:coupled-environmental-state}

The retained R1 matter--topography system is not restricted to an isolated
two-state topographic evolution.  At the anchor $z_i=2.1$, use the physical
coordinates
\begin{equation}
Y_i=
\left(
q_2,\Pi_2,
\delta_m,\delta_r,v_m,v_r
\right)_i ,
\qquad
q_2=\zeta,
\qquad
\Pi_2=2a^3G_{S,2}\dot\zeta .
\end{equation}
The same frozen linear system may be written in block form as
\begin{equation}
\begin{pmatrix}
X_2(z)\\
E(z)
\end{pmatrix}
=
\begin{pmatrix}
U_{XX}(z,z_i) & U_{XE}(z,z_i)\\
U_{EX}(z,z_i) & U_{EE}(z,z_i)
\end{pmatrix}
\begin{pmatrix}
X_2(z_i)\\
E(z_i)
\end{pmatrix},
\end{equation}
where
\begin{equation}
X_2=(q_2,\Pi_2)^T,
\qquad
E=(\delta_m,\delta_r,v_m,v_r)^T .
\end{equation}

An independent reconstruction of the frozen six-state system reproduces the
retained two-column topographic seed to $3.1\times10^{-13}$ and the stored
full canonical propagator at $z=1.5$ to $3.1\times10^{-9}$.  Setting the
initial topographic pair to zero and varying the environmental coordinates
alone gives a nonzero $2\times4$ transfer $U_{XE}$.

\begin{table}[H]
\centering
\caption{Environmental-to-topographic temporal transfer in the frozen R1
system.  The singular-value ratio depends on the coordinate normalization;
the rank and the nonzero matter-only determinant are the structural results.}
\begin{tabular}{cccc}
\toprule
$z$ & $s_1$ & $s_2$ &
$\det U_{XE}^{(\delta_m,v_m)}$\\
\midrule
1.50 & 9.65325193 & 0.00168034 & 0.0162207\\
1.00 & 7.21390025 & 0.00988894 & 0.0713376\\
0.50 & 4.13861138 & 0.02289428 & 0.0947501\\
0.00 & 3.28191222 & 0.01582572 & 0.0519383\\
\bottomrule
\end{tabular}
\end{table}

At every sampled epoch after the anchor,
\begin{equation}
\boxed{\operatorname{rank}U_{XE}=2.}
\end{equation}
Moreover, the matter-only subblock built from $(\delta_m,v_m)$ has nonzero
determinant throughout the sampled interval.  Thus both temporal components of
the later topographic state can be generated by specified environmental
density and velocity data; radiation is not required for temporal full rank.
This removes the need to postulate an isolated one-dimensional temporal
attractor for the physical $n=2$ state.

The spatial statement is separate.  The full scalar $n=2$ harmonic sector on
$S^3$ has dimension nine.  On the isotropic R1 background the same temporal
transfer acts independently on each harmonic component.  Writing
\begin{equation}
X(z)\in\mathbb R^{2\times9},
\qquad
E_i\in\mathbb R^{4\times9},
\end{equation}
one has
\begin{equation}
\boxed{
X(z)=U_{XX}(z)X_i+U_{XE}(z)E_i .
}
\end{equation}
The environment can therefore transmit an $n=2$ orientation already present
in its coefficient tensors, but isotropic R1 dynamics do not generate a
preferred axis from an isotropic state.

A one-profile reduction,
\begin{equation}
X(z)=x(z)e^T ,
\end{equation}
is consequently a special rank condition,
\begin{equation}
\boxed{
\operatorname{rank}_{\rm spatial}X(z)\le1.
}
\end{equation}
It follows automatically when the relevant environmental $n=2$ density and
velocity coefficient vectors are aligned with one common spatial profile.
Generic multi-pattern environmental data need not satisfy this condition and
can produce spatial rank two.

The resulting prediction problem is therefore not to find an isolated
self-selected topographic eigenline.  It is to specify or derive the physical
environmental state, its full $n=2$ spatial coefficients, and the justified
spatial reduction used by a given observable.

\paragraph{Claim boundary.}
This result begins at the retained R1 anchor $z_i=2.1$.  It does not derive the
anchor state from earlier cosmological initial conditions, does not determine
the absolute microscopic BHSM--R1 normalization, and does not select a unique
global spatial profile or axis.  It establishes only that, once the coupled
environmental state is specified at the anchor, the frozen R1 dynamics allow
that environment to control both temporal components of the later
topographic response without retuning.
